\documentclass[10pt,letterpaper,leqno]{amsart}
\usepackage[margin=1in]{geometry}
\usepackage[T1]{fontenc}
\usepackage{lmodern}
\usepackage{amsmath,amssymb,amsthm}
\usepackage{enumitem}
\usepackage{needspace}
\usepackage{microtype}
\usepackage[colorlinks=true,linkcolor=blue,citecolor=blue,urlcolor=blue]{hyperref}
\hypersetup{
  pdftitle={Erdos problem 597 and finite-target independence},
  pdfauthor={Will Howe, Grant Huffman, Jensen Kohlmeyer, Liam Kruer},
  pdfsubject={Ordinal graph partition relations},
  pdfkeywords={ordinal partition relations, infinite graphs, ordering property, Martin's axiom}
}
\setlength{\emergencystretch}{2em}
\raggedbottom
\clubpenalty=10000
\widowpenalty=10000
\displaywidowpenalty=10000
\numberwithin{equation}{section}
\newtheorem{theorem}{Theorem}[section]
\newtheorem{lemma}[theorem]{Lemma}
\newtheorem{proposition}[theorem]{Proposition}
\newtheorem{corollary}[theorem]{Corollary}
\theoremstyle{definition}
\newtheorem{definition}[theorem]{Definition}
\theoremstyle{remark}
\newtheorem{remark}[theorem]{Remark}
\newcommand{\Om}{\omega_1^2}
\newcommand{\Bl}{\omega_1\cdot\omega}
\newcommand{\MA}{\mathrm{MA}_{\aleph_1}}
\newcommand{\CH}{\mathrm{CH}}
\newcommand{\ZFC}{\mathrm{ZFC}}
\newcommand{\otp}{\operatorname{otp}}
\newcommand{\Gam}{\Gamma}
\newcommand{\NN}{\mathbb N}

\title[Erd\H{o}s problem 597 and finite-target independence]{Erd\H{o}s problem 597 and finite-target independence}
\author[W. Howe, G. Huffman, J. Kohlmeyer, and L. Kruer]{Will Howe, Grant Huffman, Jensen Kohlmeyer, and Liam Kruer}
\date{November 12, 2026}
\subjclass[2020]{03E02, 05C55, 03E35}
\keywords{Ordinal partition relations, infinite graphs, ordering property, Martin's axiom, continuum hypothesis}
\thanks{Electronic correspondence may be addressed to \href{mailto:howew@purdue.edu}{\texttt{howew@purdue.edu}} (W.~Howe), \href{mailto:huffma97@purdue.edu}{\texttt{huffma97@purdue.edu}} (G.~Huffman), \href{mailto:jenwin@purdue.edu}{\texttt{jenwin@purdue.edu}} (J.~Kohlmeyer), and \href{mailto:kruerl@purdue.edu}{\texttt{kruerl@purdue.edu}} (L.~Kruer).}

\begin{document}
\begin{abstract}
We study the graph-target partition relation
$\omega_1^2\to(\omega_1\cdot\omega,H)^2$ from Erd\H{o}s problem~597.
In $\ZFC$ we construct a connected bipartite graph of cardinality $\aleph_1$
and diameter at most three, containing no $K_{\omega,\omega}$, for which this
relation fails. We also construct a fixed finite $K_4$-free graph $H_0$ for
which the relation fails under the continuum hypothesis. Under Martin's
axiom at $\aleph_1$, it holds for every finite graph. Consequently the
universal finite-target assertion, and the assertion for $H_0$ itself, are
independent of $\ZFC$, relative to the consistency of $\ZFC$.
\end{abstract}
\maketitle

\section{Introduction and statements}

All graphs are simple and undirected. A copy of a graph $H$ in a graph $G$
is an injective map $f\colon V(H)\to V(G)$ satisfying
\[
 uv\in E(H)\quad\Longrightarrow\quad f(u)f(v)\in E(G).
\]
Additional edges among the image vertices are permitted. An order on the
host imposes no order requirement on this copy.

For ordinals $\alpha,\beta$ and a graph $H$, the relation
$\alpha\to(\beta,H)^2$ means that every red-blue coloring of $[\alpha]^2$
has a red homogeneous set of inherited order type $\beta$ or a blue copy
of $H$. We write
\[
 \Omega=\Om,\qquad B=\Bl.
\]
We identify $\Omega$ with $\omega_1\times\omega_1$ in row-first
lexicographic order, so $(a,b)$ has ordinal position $\omega_1a+b$.
Equivalently, the blue pairs define a graph on $\Omega$, and a red
homogeneous set is an independent set in this graph.

Erd\H{o}s asked whether $\Omega\to(B,H)^2$ holds for every graph $H$ of
cardinality at most $\aleph_1$ that contains neither $K_4$ nor
$K_{\omega,\omega}$, and whether it holds when $H$ is finite
\cite{Erdos}. Here $K_{\omega,\omega}$ denotes a complete bipartite graph
with both vertex classes countably infinite.

\begin{theorem}\label{thm:infinite}
There is a graph $H_\infty$ of cardinality $\aleph_1$ such that
\begin{enumerate}[label=\textup{(\roman*)},leftmargin=2.3em]
\item $H_\infty$ is connected and bipartite, with diameter at most three.
\item $H_\infty$ contains neither $K_4$ nor $K_{\omega,\omega}$.
\item $\Omega\not\to(B,H_\infty)^2$.
\end{enumerate}
These conclusions hold in $\ZFC$.
\end{theorem}

\begin{theorem}\label{thm:finite}
There is a fixed finite $K_4$-free graph $H_0$ such that $\CH$ implies
\begin{equation}\label{eq:negative}
 \Omega\not\to(B,H_0)^2.
\end{equation}
Under $\MA$, for every finite graph $H$,
\begin{equation}\label{eq:positive}
 \Omega\to(B,H)^2.
\end{equation}
\end{theorem}

\begin{corollary}\label{cor:independence}
Assume that $\ZFC$ is consistent. Then each of the following assertions is
independent of $\ZFC$.
\begin{enumerate}[label=\textup{(\roman*)},leftmargin=2.3em]
\item $\Omega\to(B,H)^2$ for every finite $K_4$-free graph $H$.
\item $\Omega\to(B,H_0)^2$ for the graph $H_0$ in Theorem~\ref{thm:finite}.
\end{enumerate}
\end{corollary}

The negative construction uses the nested-interval method of
Baumgartner--Hajnal \cite{BH} and Chen--Garti--Weinert \cite{CGW}.
The finite passage from an ordered obstruction to an unordered target is
an instance of the Ne\v{s}et\v{r}il--R\"odl ordering property \cite{NR}.
We give its sparse-hypergraph proof in the form needed here.
The positive partition theorem is due to Baumgartner \cite{Baumgartner}.
Related nonuniversality methods occur in Shelah \cite{Shelah}, and Li
\cite{Li} studies their application to the unrestricted graph-target
question.

\section{Uncountable rows in an ordinal block}

For $S\subseteq\Omega$ and $a<\omega_1$, put
\[
 S_a=\{b<\omega_1:(a,b)\in S\}.
\]
We call $a$ an uncountable row of $S$ when $S_a$ is uncountable. Every
uncountable subset of $\omega_1$ is unbounded and has order type $\omega_1$.

\begin{lemma}\label{lem:rows}
If $S\subseteq\Omega$ has only finitely many uncountable rows, then
$\otp(S)<B$. In particular, a subset of order type $B$ has infinitely
many uncountable rows.
\end{lemma}

\begin{proof}
Suppose the uncountable rows are $a_1<\cdots<a_k$.
Each interval of rows before $a_1$ or between successive $a_i$ contains
only countably many rows, since its indices are bounded below $\omega_1$.
All its slices are countable, so its intersection with $S$ has countable
order type. Each slice $S_{a_i}$ has order type $\omega_1$.

The part of $S$ after row $a_k$ has order type at most $\omega_1$.
Indeed, every element of that part has only countably many predecessors
there, since the preceding row indices form a countable set and every
corresponding slice is countable. A well-order in which each element has
countably many predecessors has order type at most $\omega_1$.
Adding these pieces gives
\[
 \otp(S)\le \omega_1\cdot(k+1)<\omega_1\cdot\omega.
\]
When $k=0$, the same predecessor argument applies to all of $S$.
\end{proof}

\section{An uncountable target omitted in ZFC}

Define a graph $G_*$ on $\omega_1\times\omega_1$ by
\begin{equation}\label{eq:nested}
 (a,b)\sim(r,s)
 \quad\Longleftrightarrow\quad
 a<r<s<b\ \text{ or }\ r<a<b<s.
\end{equation}

\begin{lemma}\label{lem:full-nested}
The graph $G_*$ has no independent set of order type $B$.
Every nonempty induced subgraph of $G_*$ has a vertex of at most countable
degree.
\end{lemma}

\begin{proof}
If $S$ has order type $B$, Lemma~\ref{lem:rows} gives two uncountable rows
$a<r$. Choose $s\in S_r$ with $s>r$, and then $b\in S_a$ with $b>s$.
The vertices $(a,b),(r,s)$ are adjacent by \eqref{eq:nested}.

Now let $S$ be any nonempty vertex set. Choose $(a,b)\in S$ in the least
row occupied by $S$. Every neighbor $(r,s)$ of $(a,b)$ in $S$ satisfies
$a<r<s<b$. Both its coordinates lie in the countable ordinal $b$, so
there are at most countably many such neighbors.
\end{proof}

Let $T$ be the graph on $\omega_1^{<\omega}$ in which two finite sequences
are adjacent when one is obtained from the other by appending a single
ordinal. Its cardinality is $\aleph_1$, since
\[
 |\omega_1^{<\omega}|=\sum_{n<\omega}\aleph_1^n=\aleph_1.
\]
Every vertex has $\aleph_1$ immediate extensions. The graph is a tree.
To see that it has no cycle, choose a vertex of maximum sequence length
on a purported finite cycle. Its two cycle neighbors would both be its
unique immediate predecessor, which is impossible.
Thus $T$ contains neither $K_4$ nor $K_{\omega,\omega}$, since each of
those graphs contains a finite cycle.

\begin{lemma}\label{lem:tree-omitted}
There is no copy of $T$ in $G_*$.
\end{lemma}

\begin{proof}
Under an injective edge-preserving map, the $\aleph_1$ immediate
extensions of every vertex of $T$ would give $\aleph_1$ distinct
neighbors of its image. Consequently the subgraph induced on the image
would have minimum degree at least $\aleph_1$, contrary to
Lemma~\ref{lem:full-nested}.
\end{proof}

\begin{proof}[Proof of Theorem~\ref{thm:infinite}]
The blue graph $G_*$ and its red complement already witness
$\Omega\not\to(B,T)^2$.

Partition the vertices of $T$ into $A_0,A_1$ according to the parity of
sequence length. Add vertices $h_0,h_1$, place $h_i$ in the side $A_i$,
and join $h_i$ to every vertex in the opposite side, including $h_{1-i}$.
Call the resulting graph $H_\infty$.
It is bipartite and has cardinality $\aleph_1$.
Two old vertices in the same side are joined through the opposite hub
by a path of length two. Two old vertices in opposite sides are joined
through the two hubs by a path of length three. A pair involving a hub
has distance at most two. Hence $H_\infty$ is connected and has diameter
at most three. Bipartiteness implies that it is $K_4$-free.

If $H_\infty$ contained a $K_{\omega,\omega}$, removing the two hubs
would leave two infinite vertex classes with every cross-edge present
in $T$, a contradiction. Finally, any copy of $H_\infty$ in $G_*$ would
restrict to a copy of $T$. Lemmas~\ref{lem:full-nested} and
\ref{lem:tree-omitted} therefore give the required negative partition
relation for $H_\infty$.
\end{proof}

\section{A sparse nested graph under CH}

Assume $\CH$. For every $t<\omega_1$, fix an injection
\[
 e_t\colon t\longrightarrow\NN.
\]
A \emph{pattern} is a countable set $P$ of vertices $(r,s)$ with $r<s$
such that $P$ occupies infinitely many rows and contains infinitely many
vertices in each row that it occupies. Since
\[
 (\aleph_1)^{\aleph_0}
 =(2^{\aleph_0})^{\aleph_0}
 =2^{\aleph_0}=\aleph_1,
\]
we may enumerate all patterns as $(P_\nu)_{\nu<\omega_1}$, allowing
repetitions.

For $x=(a,b)$ we construct an outgoing neighborhood $\Gam(x)$.
The construction proceeds by recursion on $b$.
If $a\ge b$, set $\Gam(x)=\varnothing$.
If $a<b$, call an index $\nu$ \emph{eligible for $x$} when
\begin{equation}\label{eq:eligible}
 \nu<b
 \quad\text{and}\quad
 a<r<s<b\quad\text{for every }(r,s)\in P_\nu.
\end{equation}
The eligible indices form a countable set. List them without repetition
in a sequence of length at most $\omega$.
At each step, choose one point $(r,s)$ from the current eligible pattern
subject to the following rules.
\begin{enumerate}[label=\textup{(\roman*)},leftmargin=2.3em]
\item Its row has not previously been selected into $\Gam(x)$.
\item For every previously selected row $t>r$,
\begin{equation}\label{eq:selection-rank}
 e_t(r)>e_t(a),
 \qquad
 e_t(r)>e_t(u)
 \quad\text{for every previously selected row }u<t.
\end{equation}
\item For every previously selected point $(u,v)$ with $u<r$,
\begin{equation}\label{eq:selection-avoid}
 (r,s)\notin\Gam(u,v).
\end{equation}
\end{enumerate}
Let $G$ be the undirected graph whose edges are the pairs $xy$ with
$y\in\Gam(x)$.

\begin{lemma}\label{lem:construction}
Every selection step is possible. The resulting graph satisfies
\begin{enumerate}[label=\textup{(\roman*)},leftmargin=2.3em]
\item if $(r,s)\in\Gam(a,b)$, then $a<r<s<b$.
\item $\Gam(a,b)$ contains at most one vertex in each row.
\item $\Gam(a,b)$ meets every pattern eligible for $(a,b)$.
\end{enumerate}
\end{lemma}

\begin{proof}
There are only finitely many previous selections at a given step.
Rule~\textup{(i)} excludes finitely many rows.
For a fixed previously selected row $t$, let $M_t$ be the maximum of
$e_t(a)$ and the finitely many values $e_t(u)$ with $u<t$ previously
selected. Injectivity of $e_t$ implies that
\[
 \{r<t:e_t(r)\le M_t\}
\]
is finite. Thus rule~\textup{(ii)} also excludes only finitely many rows.
The current pattern has infinitely many rows, so an admissible row $r$
can be chosen.

At this row, each set $\Gam(u,v)$ occurring in
rule~\textup{(iii)} is already defined, since $v<b$ by
\eqref{eq:eligible}. By induction it contains at most one point in
row $r$. Consequently rule~\textup{(iii)} excludes finitely many points
from an infinite slice of the pattern. A permissible $(r,s)$ exists.
The three asserted properties follow from eligibility, rule~\textup{(i)},
and the processing of every eligible index.
\end{proof}

If vertices $x,y$ lie in increasing rows, then
\begin{equation}\label{eq:orientation}
 xy\in E(G)\quad\Longleftrightarrow\quad y\in\Gam(x).
\end{equation}
In particular, $G$ has no same-row edges.

\Needspace{6\baselineskip}
\begin{lemma}\label{lem:no-red}
The graph $G$ has no independent set of order type $B$.
\end{lemma}

\begin{proof}
Let $S\subseteq\Omega$ have order type $B$.
By Lemma~\ref{lem:rows}, choose an uncountable row $a$ of $S$ with
infinitely many uncountable rows of $S$ above it.
Choose a countably infinite collection of those later rows, and in each
chosen row $r$ choose infinitely many points $(r,s)\in S$ with $s>r$.
Their union is a pattern $P_\nu\subseteq S$.

The coordinates of $P_\nu$, together with $\nu$, are bounded below
$\omega_1$. Since $S_a$ is unbounded, choose $b\in S_a$ strictly above
all these ordinals. Then $P_\nu$ is eligible for $(a,b)$.
Lemma~\ref{lem:construction} supplies a point of
$\Gam(a,b)\cap P_\nu$, giving an edge of $G$ with both ends in $S$.
\end{proof}

\section{The ordered six-vertex obstruction}

For $y\in\Gam(x)$, let $\tau_x(y)$ be its selection index in the
construction of $\Gam(x)$. The uniqueness of its selected row makes
this index well defined.

\begin{lemma}[Selection reversal]\label{lem:reversal}
Suppose $u,v\in\Gam(x)$, the row of $u$ is smaller than the row of $v$,
and $uv\in E(G)$. Then $\tau_x(v)<\tau_x(u)$.
\end{lemma}

\begin{proof}
Equation~\eqref{eq:orientation} gives $v\in\Gam(u)$.
If $u$ had been selected before $v$ in $\Gam(x)$,
rule~\eqref{eq:selection-avoid} would have excluded $v$.
The two selection indices are distinct, so the stated inequality follows.
\end{proof}

\begin{lemma}[Triangle comparison]\label{lem:triangle}
If $u,v,w$ form a triangle of $G$ in rows $r_u<r_v<r_w$, then
\begin{equation}\label{eq:triangle}
 e_{r_w}(r_u)<e_{r_w}(r_v).
\end{equation}
\end{lemma}

\begin{proof}
Both $v,w$ lie in $\Gam(u)$.
Lemma~\ref{lem:reversal} places $w$ before $v$ in that neighborhood.
Rule~\eqref{eq:selection-rank}, applied when $v$ is selected and with
$t=r_w$, gives \eqref{eq:triangle}.
\end{proof}

Let $F$ be the complete tripartite graph with parts
\[
 \{x,p\},\qquad \{y,z\},\qquad \{q,w\},
\]
equipped with the order
\begin{equation}\label{eq:order-F}
 x<y<p<q<z<w.
\end{equation}

\begin{proposition}\label{prop:ordered-obstruction}
There is no edge-preserving copy of $F$ in $G$ whose six image rows are
strictly increasing in the order \eqref{eq:order-F}.
\end{proposition}

\begin{proof}
Suppose such a copy exists, and retain the labels of $F$ for its image
vertices. Write $r_v$ for the row of an image vertex $v$.
The vertices $y,q,z,w$ lie in $\Gam(x)$.
Applying Lemma~\ref{lem:reversal} to the edges $zw,qz,yq$ gives
\[
 \tau_x(w)<\tau_x(z)<\tau_x(q)<\tau_x(y).
\]
At the selection of $y$, the row $r_w$ has already been selected, as has
$r_z<r_w$. Rule~\eqref{eq:selection-rank} therefore gives
\begin{equation}\label{eq:rank-backward}
 e_{r_w}(r_z)<e_{r_w}(r_y).
\end{equation}
The triples $y,p,w$ and $p,z,w$ are triangles of $F$.
Lemma~\ref{lem:triangle} gives
\[
 e_{r_w}(r_y)<e_{r_w}(r_p)<e_{r_w}(r_z),
\]
contradicting \eqref{eq:rank-backward}.
\end{proof}

\begin{proposition}\label{prop:host-K4}
The graph $G$ is $K_4$-free.
\end{proposition}

\begin{proof}
A $K_4$ would occupy four distinct rows $r_a<r_b<r_c<r_d$.
Its last three vertices lie in $\Gam(a)$.
Selection reversal gives
$\tau_a(d)<\tau_a(c)<\tau_a(b)$.
At the selection of $b$, rule~\eqref{eq:selection-rank} implies
$e_{r_d}(r_c)<e_{r_d}(r_b)$.
The triangle $b,c,d$ gives the opposite strict inequality by
Lemma~\ref{lem:triangle}.
\end{proof}

\Needspace{12\baselineskip}
\section{A finite unordered target}

The following specialization of the Ne\v{s}et\v{r}il--R\"odl ordering
property \cite{NR} removes the prescribed order on $F$.

\begin{lemma}\label{lem:ordering-property}
There is a finite $K_4$-free graph $H_0$ such that every linear order on
$V(H_0)$ contains an order-preserving, edge-preserving copy of the
ordered graph $F$ from \eqref{eq:order-F}.
\end{lemma}

\begin{proof}
For a sufficiently large integer $n$, choose a random $6$-uniform
hypergraph on $[n]$, including each six-element set independently with
probability $p=n^{-19/4}$. Let $X$ count its hyperedges. Let $Y$ count
pairs of distinct hyperedges with an intersection of size at least two.
Let $Z$ count triples of distinct hyperedges $E_1,E_2,E_3$ admitting
distinct vertices
\[
 v_{12}\in E_1\cap E_2,\qquad
 v_{23}\in E_2\cap E_3,\qquad
 v_{31}\in E_3\cap E_1.
\]
There are at most $n^{10}$ possible pairs counted by $Y$.
Indeed, choose two common vertices and at most four further vertices
for each of the two hyperedges. There are at most $n^{15}$ possible
triples counted by $Z$, by choosing the three intersection witnesses
and at most four further vertices for each hyperedge. Thus
\begin{align*}
 \mathbb E(X-Y-Z)
 &\ge \binom n6 p-n^{10}p^2-n^{15}p^3\\
 &=\left(\frac1{720}+o(1)\right)n^{5/4}
       -n^{1/2}-n^{3/4}.
\end{align*}
There is an absolute $c>0$ such that, for all sufficiently large $n$,
some outcome satisfies $X-Y-Z\ge c n^{5/4}$.
From this outcome, delete one hyperedge from each originally counted
bad pair or triple. Repeated deletions of the same hyperedge are counted
once. The resulting hypergraph $\mathcal H$ has
\begin{equation}\label{eq:many-hyperedges}
 m=|\mathcal H|\ge c n^{5/4},
\end{equation}
distinct hyperedges intersect in at most one vertex, and no triple of
hyperedges has three distinct intersection witnesses as above.

Independently for each $E\in\mathcal H$, choose a uniformly random
bijection from the six labels of $F$ onto $E$, and place the corresponding
edges of $F$ on $E$. Let $H_0$ be the union of all these graph edges.
Every graph edge has a unique supporting hyperedge, since two
hyperedges meet in at most one vertex.

Every triangle of $H_0$ lies in one supporting hyperedge.
If two of its edges have the same support $E$, all three triangle
vertices belong to $E$, and the support of its third edge must also be
$E$. If the three supports are distinct, the three triangle vertices
are distinct intersection witnesses, contrary to the construction of
$\mathcal H$.
If $H_0$ contained a $K_4$, two triangles sharing an edge would have the
same supporting hyperedge. All four vertices would lie in that
hyperedge, and uniqueness of support would place all six edges in its
copy of $F$. This contradicts the fact that $F=K_{2,2,2}$ is $K_4$-free.
Therefore every choice of the label bijections produces a $K_4$-free
graph $H_0$.

Fix a linear order on $[n]$. For any $E\in\mathcal H$, the probability
that its labels occur in precisely the order \eqref{eq:order-F} is
$1/6!$. These events are independent over the hyperedges, since their
label bijections were chosen independently. The probability that at
least one of the $n!$ linear orders has no successful hyperedge is at most
\begin{equation}\label{eq:union-bound}
 n!\left(1-\frac1{6!}\right)^m
 \le \exp\left(n\log n-\frac{m}{720}\right)<1
\end{equation}
for sufficiently large $n$, by \eqref{eq:many-hyperedges}.
Some choice of the label bijections therefore works for every linear
order on $[n]$. Fix that choice to obtain $H_0$.
\end{proof}

\begin{proposition}\label{prop:no-finite-copy}
Under $\CH$, the graph $G$ from Section~4 contains no copy of $H_0$.
\end{proposition}

\begin{proof}
Suppose $f\colon V(H_0)\to V(G)$ is an injective edge-preserving map.
Order $V(H_0)$ by the rows of its images under $f$, breaking ties by
any fixed linear order. Lemma~\ref{lem:ordering-property} gives an
ordered copy of $F$ for this linear order.
Its six image rows are nondecreasing in the order
$x,y,p,q,z,w$. Each consecutive pair in this list is an edge of $F$.
Since $G$ has no same-row edges, each consecutive pair of image rows
is strictly increasing. Proposition~\ref{prop:ordered-obstruction}
now gives a contradiction.
\end{proof}

Color the edges of $G$ blue and all other pairs red.
Lemma~\ref{lem:no-red} and Proposition~\ref{prop:no-finite-copy}
show that this coloring has neither a red set of order type $B$ nor a
blue copy of $H_0$. The finite graph $H_0$ was fixed in
Lemma~\ref{lem:ordering-property}, before the assumption of $\CH$.
This proves \eqref{eq:negative}.

\section{The positive relation and independence}

Recall that $\MA$ states that for every nonempty partially ordered set satisfying
the countable chain condition, every family of at most $\aleph_1$ dense
subsets is met by a filter. We use the following partition theorem.

\begin{theorem}[Baumgartner {\cite[Theorem~3.5(a)]{Baumgartner}}]\label{thm:Baumgartner}
Assume $\MA$. Then
\[
 \Bl\to(\Bl,n)^2
 \qquad\text{for every finite }n.
\]
\end{theorem}

Here the second alternative is a blue complete graph on $n$ vertices.

\begin{proof}[Proof of \eqref{eq:positive} and Theorem~\ref{thm:finite}]
Let $H$ be a finite graph with $n$ vertices. For $n=0$, every host contains
the empty copy, so the assertion is immediate.
For $n>0$, restrict an arbitrary coloring of $[\Omega]^2$ to its initial
segment of type $B$. Theorem~\ref{thm:Baumgartner} gives a red set of
order type $B$ or a blue $K_n$. Every $n$-vertex graph has an
edge-preserving copy in $K_n$, proving \eqref{eq:positive}.
Together with the construction in Section~6, this proves
Theorem~\ref{thm:finite}.
\end{proof}

\begin{proof}[Proof of Corollary~\ref{cor:independence}]
Let $\mathsf A$ be the assertion that $\Omega\to(B,H)^2$ for every
finite $K_4$-free graph $H$, and let $\mathsf A_0$ be
$\Omega\to(B,H_0)^2$.
Theorem~\ref{thm:finite} proves
\[
 \CH\ \Longrightarrow\ \neg\mathsf A\wedge\neg\mathsf A_0,
 \qquad
 \MA\ \Longrightarrow\ \mathsf A\wedge\mathsf A_0.
\]
G\"odel's relative consistency theorem gives a model of
$\ZFC+\CH$ whenever $\ZFC$ is consistent \cite{Godel}.
The relative consistency theorem of Solovay and Tennenbaum gives a model
of $\ZFC+\MA$ under the same assumption \cite{ST}.
Thus each of $\mathsf A$ and $\mathsf A_0$, as well as each of their
negations, is consistent with $\ZFC$ relative to its consistency.
\end{proof}

\section*{Acknowledgments}
The authors acknowledge assistance from OpenAI's ChatGPT in mathematical
exploration, proof checking, and preparation of the exposition.

\begin{thebibliography}{9}
\raggedright

\bibitem{Baumgartner}
J.~E. Baumgartner,
\emph{Remarks on partition ordinals},
in \emph{Set Theory and its Applications} (J.~Stepr\={a}ns and
S.~Watson, eds.), Lecture Notes in Mathematics \textbf{1401},
Springer, Berlin, 1989, 5--17.
\href{https://doi.org/10.1007/BFb0097328}{doi:10.1007/BFb0097328}.

\bibitem{BH}
J.~E. Baumgartner and A.~Hajnal,
\emph{A remark on partition relations for infinite ordinals with an
application to finite combinatorics},
in \emph{Logic and Combinatorics} (S.~G. Simpson, ed.),
Contemporary Mathematics \textbf{65}, American Mathematical Society,
Providence, RI, 1987, 157--167.
\href{https://doi.org/10.1090/conm/065/891246}{doi:10.1090/conm/065/891246}.

\bibitem{CGW}
W.~Chen, S.~Garti, and T.~Weinert,
\emph{Cardinal characteristics of the continuum and partitions},
Israel Journal of Mathematics \textbf{235} (2020), 13--38.
\href{https://doi.org/10.1007/s11856-019-1942-y}{doi:10.1007/s11856-019-1942-y}.

\bibitem{Erdos}
P.~Erd\H{o}s,
\emph{Some problems on finite and infinite graphs},
in \emph{Logic and Combinatorics} (S.~G. Simpson, ed.),
Contemporary Mathematics \textbf{65}, American Mathematical Society,
Providence, RI, 1987, 223--228.

\bibitem{Godel}
K.~G\"odel,
\emph{The Consistency of the Axiom of Choice and of the Generalized
Continuum-Hypothesis with the Axioms of Set Theory},
Annals of Mathematics Studies \textbf{3},
Princeton University Press, Princeton, NJ, 1940.

\bibitem{Li}
A.~C. Li,
\emph{A nonuniversality obstruction to an ordinal graph partition relation},
preprint, version~1.1, 2026.
\href{https://doi.org/10.5281/zenodo.22671917}{doi:10.5281/zenodo.22671917}.

\bibitem{NR}
J.~Ne\v{s}et\v{r}il and V.~R\"odl,
\emph{On a probabilistic graph-theoretical method},
Proceedings of the American Mathematical Society \textbf{72} (1978),
417--421.
\href{https://doi.org/10.1090/S0002-9939-1978-0507350-7}{doi:10.1090/S0002-9939-1978-0507350-7}.

\bibitem{Shelah}
S.~Shelah,
\emph{Universality among graphs omitting a complete bipartite graph},
Combinatorica \textbf{32} (2012), 325--362.
\href{https://doi.org/10.1007/s00493-012-2033-4}{doi:10.1007/s00493-012-2033-4}.

\bibitem{ST}
R.~M. Solovay and S.~Tennenbaum,
\emph{Iterated Cohen extensions and Souslin's problem},
Annals of Mathematics (2) \textbf{94} (1971), 201--245.
\href{https://doi.org/10.2307/1970860}{doi:10.2307/1970860}.

\end{thebibliography}
\end{document}
